\documentclass[11pt,a4paper]{article}

%================================================
% ENCODAGE ET LANGUE
%================================================
\usepackage[utf8]{inputenc}
\usepackage[T1]{fontenc}
\usepackage[french]{babel}

%================================================
% MISE EN PAGE
%================================================
\usepackage{geometry}
\geometry{margin=1.55cm}

%================================================
% PACKAGE GTOLI
%================================================
\usepackage{../styles/gtoli}

%================================================
% INFORMATIONS DU DOCUMENT
%================================================
\title{\textbf{Mathématiques}\\[0.2cm]
\large Titre du chapitre}

\author{GToli}
\date{\today}

%================================================
% DOCUMENT
%================================================
\begin{document}

\maketitle

%================================================
\section{Notion concernée}
%================================================

\begin{cadreobjectif}
Dans ce chapitre, l’élève doit être capable de :
\begin{itemize}
    \item identifier la notion mathématique étudiée ;
    \item comprendre les définitions et propriétés associées ;
    \item appliquer une méthode de résolution structurée ;
    \item justifier les étapes du raisonnement ;
    \item rédiger une solution mathématique correcte ;
    \item résoudre des exercices progressifs et des situations de transfert.
\end{itemize}
\end{cadreobjectif}

%================================================
\section{Prérequis}
%================================================

\begin{cadregris}{Notions nécessaires}
Avant de commencer, il faut maîtriser :
\begin{itemize}
    \item les notations mathématiques utiles ;
    \item les règles de calcul de base ;
    \item les propriétés déjà vues ;
    \item les conventions propres au chapitre ;
    \item les techniques élémentaires nécessaires à la résolution.
\end{itemize}
\end{cadregris}

%================================================
\section{Rappel théorique}
%================================================

\begin{cadretheorie}
Cette section contient :
\begin{itemize}
    \item les définitions ;
    \item les propriétés ;
    \item les théorèmes ;
    \item les formules ;
    \item les conditions d’application ;
    \item les notations importantes.
\end{itemize}

\medskip

\textbf{Point essentiel :} une propriété mathématique ne doit pas seulement être connue ; elle doit être utilisée dans un raisonnement correctement justifié.
\end{cadretheorie}

%================================================
\section{Démonstration ou justification}
%================================================

\begin{preuve}
Cette section peut contenir :
\begin{itemize}
    \item une démonstration ;
    \item une justification d’une formule ;
    \item une explication du lien entre deux propriétés ;
    \item une mise en évidence des conditions d’application.
\end{itemize}
\end{preuve}

%================================================
\section{Méthode générale de résolution}
%================================================

\begin{cadremethode}
Pour résoudre un exercice mathématique :

\begin{enumerate}
    \item lire précisément l’énoncé ;
    \item identifier la nature de la question ;
    \item repérer les données utiles ;
    \item choisir la propriété ou la méthode adaptée ;
    \item vérifier les conditions d’application ;
    \item effectuer les calculs de manière structurée ;
    \item justifier les étapes importantes ;
    \item contrôler la cohérence du résultat ;
    \item rédiger une conclusion claire.
\end{enumerate}
\end{cadremethode}

%================================================
\section{Tableau de synthèse}
%================================================

\begin{cadregris}{Synthèse des méthodes}
\renewcommand{\arraystretch}{1.35}
\begin{tabularx}{\textwidth}{>{\bfseries}p{3.5cm}X X}
\toprule
Type de question & Indice dans l’énoncé & Méthode à mobiliser \\
\midrule
Calcul direct & Données explicites & Appliquer la propriété \\
Justification & « Montrer que », « justifier » & Rédiger un raisonnement \\
Étude de signe & Expression factorisée ou à factoriser & Tableau de signes \\
Variation & Présence d’une fonction & Étudier la dérivée \\
Problème contextualisé & Situation réelle ou modélisée & Traduire mathématiquement \\
\bottomrule
\end{tabularx}
\end{cadregris}

%================================================
\section{Exemple modèle}
%================================================

\begin{cadreexemple}
\textbf{Énoncé :}

Insérer ici un exemple représentatif du chapitre.

\medskip

\textbf{Analyse :}
\begin{itemize}
    \item Quelle est la nature de la question ?
    \item Quelles données sont disponibles ?
    \item Quelle propriété faut-il utiliser ?
    \item Quelle méthode est la plus efficace ?
\end{itemize}

\medskip

\textbf{Résolution commentée :}

Présenter les étapes avec justification.
\end{cadreexemple}

%================================================
\section{Protocole spécifique}
%================================================

\begin{cadremethode}
Cette section peut être utilisée pour un protocole propre au chapitre.

\medskip

Exemples :
\begin{itemize}
    \item protocole d’étude de fonction ;
    \item protocole de résolution d’équation ;
    \item protocole de probabilité ;
    \item protocole de démonstration ;
    \item protocole de lecture graphique.
\end{itemize}
\end{cadremethode}

%================================================
\section{Erreurs fréquentes}
%================================================

\begin{cadreattention}
Les erreurs fréquentes en mathématiques sont :
\begin{itemize}
    \item appliquer une formule sans vérifier ses conditions ;
    \item confondre deux propriétés proches ;
    \item oublier une condition d’existence ;
    \item commettre une erreur de signe ;
    \item négliger la rédaction ;
    \item donner un résultat sans justification ;
    \item ne pas interpréter le résultat dans le contexte.
\end{itemize}
\end{cadreattention}

%================================================
\section{Exercice d’application}
%================================================

\begin{cadreenonce}
\textbf{Exercice 1 :}

Insérer ici un exercice d’application directe.
\end{cadreenonce}

\begin{cadreaide}
Pour commencer :
\begin{itemize}
    \item identifier la notion principale ;
    \item repérer les données utiles ;
    \item choisir la bonne propriété ;
    \item organiser les calculs ;
    \item rédiger la conclusion.
\end{itemize}
\end{cadreaide}

%================================================
\section{Solution détaillée}
%================================================

\begin{cadresolution}
La solution doit comprendre :
\begin{enumerate}
    \item l’analyse de l’énoncé ;
    \item le choix de la méthode ;
    \item les calculs ;
    \item les justifications ;
    \item la conclusion.
\end{enumerate}
\end{cadresolution}

%================================================
\section{Exercices progressifs}
%================================================

\begin{cadreactivite}
\begin{enumerate}
    \item Application directe.
    \item Application avec adaptation.
    \item Exercice à plusieurs étapes.
    \item Exercice de transfert.
    \item Exercice de synthèse.
    \item Exercice type évaluation.
\end{enumerate}
\end{cadreactivite}

%================================================
\section{Exercice de transfert}
%================================================

\begin{cadreenonce}
\textbf{Exercice de transfert :}

Insérer ici une situation nécessitant une adaptation de la méthode.
\end{cadreenonce}

%================================================
\section{Synthèse finale}
%================================================

\begin{cadresynthese}
À retenir :

\begin{itemize}
    \item notion centrale : \dotfill
    \item propriété importante : \dotfill
    \item méthode principale : \dotfill
    \item condition d’application : \dotfill
    \item erreur fréquente : \dotfill
    \item stratégie de vérification : \dotfill
\end{itemize}
\end{cadresynthese}

%================================================
\section{Auto-évaluation}
%================================================

\begin{cadregris}{Je vérifie ma maîtrise}
\begin{itemize}
    \item[$\square$] Je connais les définitions.
    \item[$\square$] Je sais choisir la méthode.
    \item[$\square$] Je sais appliquer la propriété.
    \item[$\square$] Je sais justifier les étapes.
    \item[$\square$] Je sais rédiger une solution.
    \item[$\square$] Je peux résoudre une variante.
\end{itemize}
\end{cadregris}

\end{document}